\begin{figure}[ht]
\centering
\begin{tikzpicture}[x=1cm, y=1cm, font=\small]

% Error vs term count for raw pi/4 partial sums, single Aitken, Euler.
% Data from data/acceleration_compare.csv.
% x: term count N = 2..14 mapped x_cm = (N-2)*0.65.
% y: log10(error). range 10^0 (top) to 10^-6 (bottom). y_cm = 6 + log10(err).

\draw[->, thick] (-0.3, 0) -- (8.7, 0) node[right] {$N$ (terms used)};
\draw[->, thick] (0, -0.3) -- (0, 7.0) node[above] {error (log scale)};

\foreach \n in {2,4,6,8,10,12,14} {
  \pgfmathsetmacro\xc{(\n-2)*0.65}
  \draw (\xc, -0.1) -- (\xc, 0.1) node[below, yshift=-0.12cm] {\n};
}
\foreach \p in {0,-1,-2,-3,-4,-5,-6} {
  \pgfmathsetmacro\yc{6+\p}
  \draw (-0.1, \yc) -- (0.1, \yc) node[left, xshift=-0.1cm] {$10^{\p}$};
}

% Raw partial sums of pi/4: err ~ 0.12 down to 0.018 (log10 -0.92 .. -1.75)
% y_cm = 6 + log10(err): N=2 -0.93->5.07; N=4 -1.21->4.79; N=6 -1.38->4.62;
% N=8 -1.51->4.49; N=10 -1.60->4.40; N=12 -1.68->4.32; N=14 -1.75->4.25
\draw[thick, engblack, mark=*, mark size=1pt] plot coordinates {
  (0,5.07) (1.3,4.79) (2.6,4.62) (3.9,4.49) (5.2,4.40) (6.5,4.32) (7.8,4.25) };
\node[engblack, anchor=west] at (7.85, 4.25) {raw};

% Single Aitken: N=4 2.065e-3 ->3.31; N=6 4.78e-4 ->2.68; N=8 1.78e-4 ->2.25;
% N=10 8.45e-5 ->1.93; N=12 4.65e-5 ->1.67; N=14 2.82e-5 ->1.45
\draw[thick, engblue, mark=square*, mark size=1pt] plot coordinates {
  (1.3,3.31) (2.6,2.68) (3.9,2.25) (5.2,1.93) (6.5,1.67) (7.8,1.45) };
\node[engblue, anchor=west] at (7.85, 1.45) {Aitken};

% Euler transform: N=4 2.35e-2 ->4.37; N=6 5.02e-3 ->3.70; N=8 1.12e-3 ->3.05;
% N=10 2.54e-4 ->2.40; N=12 5.86e-5 ->1.77; N=14 1.37e-5 ->1.14
\draw[thick, engred, mark=triangle*, mark size=1pt] plot coordinates {
  (1.3,4.37) (2.6,3.70) (3.9,3.05) (5.2,2.40) (6.5,1.77) (7.8,1.14) };
\node[engred, anchor=west] at (7.85, 1.14) {Euler};

\end{tikzpicture}
\caption[Convergence acceleration of the Leibniz series for $\pi/4$]{Error
of the raw Leibniz partial sums for $\pi/4 = 1 - 1/3 + 1/5 - \cdots$
against the error of the single-pass Aitken-$\Delta^{2}$ and
Euler-transformed estimates, as a function of the number of raw terms
used. The raw series barely improves, tracing the slow $O(1/N)$ harmonic
rate. The accelerated estimates fall two to three decades below it on the
same term budget, and iterating the Aitken transform on its own output (not
shown) reaches six correct digits from the first eight raw terms. The
convergence test settles that the limit exists; the acceleration transform
delivers it cheaply.}
\label{fig:vol02:ch04:acceleration}
\end{figure}
